% A \label in vertical mode – straight after \begin{theorem}, \section and \item, and a theorem straight after a display – leaves the spacing as it is without the label.
\documentclass{article}
\usepackage{amsmath,amsthm}
\newtheorem{theorem}{Theorem}
\newtheorem{definition}{Definition}
% one page with nothing but the text: checked against plain lualatex too (cases.ts)
\pagestyle{empty}
\begin{document}
\section{Introduction}\label{s:intro}
A paragraph under a labelled section, long enough to break onto more than one
line of the column, and with a descender to end on: quickly.
\begin{theorem}[Labelled after a paragraph]\label{t:para}
A theorem whose label comes straight after its \verb|\begin|, before its
text, as it does in most papers that label their results.
\end{theorem}
A paragraph that ends in a display, set on its own line, so that the next
environment starts straight below it:
\begin{equation}
  f(y) = \int_0^1 g(x)\,dx.
\end{equation}
\begin{definition}\label{d:display}
A definition whose label comes first, straight after the display above, with
no paragraph between them.
\end{definition}
Another display, then a blank line before a labelled environment:
\[
  \sum_{k=1}^{n} k = \frac{n(n+1)}{2}.
\]

\begin{theorem}\label{t:blank}
A theorem after a display and a blank line, labelled first as well.
\end{theorem}
\section{Lists}\label{s:lists}
\begin{itemize}
  \item\label{i:one} A first item with its label before its text, long enough
        to break onto a second line when the column is narrow, quickly.
  \item\label{i:two} A second item, labelled the same way.
\end{itemize}
A closing paragraph, so that the last label is followed by text set from the
depth of the line above it.
\end{document}
